\documentclass{article}
\usepackage{amsmath,amssymb,amsthm}
\usepackage{algorithm}
\usepackage{algorithmic}
\usepackage{enumitem}
\usepackage{hyperref}
\newtheorem{theorem}{Theorem}
\newtheorem{lemma}{Lemma}
\newtheorem{assumption}{Assumption}
\newtheorem{definition}{Definition}
\newcommand{\norm}[1]{\left\lVert#1\right\rVert}
\newcommand{\R}{\mathbb{R}}
\newcommand{\E}{\mathbb{E}}

\begin{document}

\section*{Trust Region Method with Approximate Subproblem Solver (TR‑AP)}

\section*{Mathematical Formulation}
Let $f:\R^{n}\rightarrow\R$ be the objective function.  
At iteration $k$ we have the current iterate $x_{k}\in\R^{n}$ and a trust‑region radius
$\Delta_{k}>0$.  
Define the (quadratic) model
\[
m_{k}(p)\;:=\;f(x_{k})+g_{k}^{\top}p+\tfrac12 p^{\top}B_{k}p,
\qquad 
g_{k}\;:=\;\nabla f(x_{k}),\; B_{k}\approx\nabla^{2}f(x_{k}).
\]
The (exact) trust‑region subproblem is
\[
\min_{p\in\R^{n}}\; m_{k}(p)\quad\text{s.t.}\quad \norm{p}\le \Delta_{k}.
\tag{TR}
\]
In TR‑AP we solve \eqref{TR} only approximately.  
A standard inexpensive approximant is the \emph{Cauchy point}
\[
p_{k}\;:=\; -\frac{\Delta_{k}}{\norm{g_{k}}}\,g_{k},
\tag{1}
\]
which is feasible and yields a guaranteed amount of model reduction (see Lemma~1).  

Define the \emph{actual reduction} and \emph{predicted reduction}
\[
\text{ared}_{k}\;:=\;f(x_{k})-f(x_{k}+p_{k}),\qquad
\text{pred}_{k}\;:=\;m_{k}(0)-m_{k}(p_{k}).
\]
The ratio
\[
\rho_{k}\;:=\;\frac{\text{ared}_{k}}{\text{pred}_{k}}
\tag{2}
\]
determines whether the step is accepted and how the radius is updated:
\[
\begin{cases}
\text{if }\rho_{k}\ge \eta_{1} \;(0<\eta_{1}<1), & 
   x_{k+1}=x_{k}+p_{k},\;\Delta_{k+1}=
   \min\{\gamma_{\text{inc}}\Delta_{k},\Delta_{\max}\},\\[0.2em]
\text{if }\rho_{k}<\eta_{1}, &
   x_{k+1}=x_{k},\;\Delta_{k+1}= \gamma_{\text{dec}}\Delta_{k},
\end{cases}
\tag{3}
\]
with parameters $\gamma_{\text{inc}}>1$, $0<\gamma_{\text{dec}}<1$ and a maximum radius $\Delta_{\max}>0$.

\section*{Pseudocode}
\begin{algorithm}[H]
\caption{TR‑AP}
\begin{algorithmic}[1]
\STATE \textbf{Input:} $x_{0}\in\R^{n}$, $\Delta_{0}>0$, 
$\eta_{1}\in(0,1)$, $\gamma_{\text{inc}}>1$, $0<\gamma_{\text{dec}}<1$, $\Delta_{\max}>0$.
\FOR{$k=0,1,2,\dots$}
    \STATE Compute $g_{k}=\nabla f(x_{k})$, choose $B_{k}$ (e.g. $B_{k}= \nabla^{2}f(x_{k})$ or $B_{k}=I$).
    \STATE Form the Cauchy step \eqref{1}: $p_{k}= -\dfrac{\Delta_{k}}{\norm{g_{k}}}g_{k}$.
    \STATE Evaluate $f(x_{k}+p_{k})$ and compute $\rho_{k}$ via \eqref{2}.
    \IF{$\rho_{k}\ge \eta_{1}$}
        \STATE $x_{k+1}\gets x_{k}+p_{k}$.
        \STATE $\Delta_{k+1}\gets\min\{\gamma_{\text{inc}}\Delta_{k},\Delta_{\max}\}$.
    \ELSE
        \STATE $x_{k+1}\gets x_{k}$.
        \STATE $\Delta_{k+1}\gets\gamma_{\text{dec}}\Delta_{k}$.
    \ENDIF
    \STATE \textbf{Termination test:} if $\norm{g_{k}}\le\varepsilon$ stop.
\ENDFOR
\end{algorithmic}
\end{algorithm}

\section*{Dry Run Example}
Consider the one‑dimensional problem
\[
f(w)=(w-3)^{2},\qquad w\in\R.
\]
Then $g(w)=\nabla f(w)=2(w-3)$ and $B(w)=\nabla^{2}f(w)=2$.  
Choose $w_{0}=0$, initial radius $\Delta_{0}=1$, and parameters
$\eta_{1}=0.1$, $\gamma_{\text{inc}}=2$, $\gamma_{\text{dec}}=0.5$.
We apply the Cauchy step \eqref{1} for three iterations.

\begin{center}
\begin{tabular}{c|c|c|c|c}
$k$ & $w_{k}$ & $g_{k}=2(w_{k}-3)$ & $p_{k}= -\dfrac{\Delta_{k}}{|g_{k}|}g_{k}$ & $f(w_{k})$\\ \hline
0 & $0$ & $-6$ & $- \frac{1}{6}(-6)= 1$ & $(0-3)^{2}=9$\\
1 & $1$ & $-4$ & $- \frac{1}{4}(-4)= 1$ & $(1-3)^{2}=4$\\
2 & $2$ & $-2$ & $- \frac{1}{2}(-2)= 1$ & $(2-3)^{2}=1$\\
3 & $3$ & $0$  & $0$ & $(3-3)^{2}=0$
\end{tabular}
\end{center}
At each iteration $\rho_{k}=1$ (the model is exact for a quadratic), so the step is accepted and the radius is enlarged. After three steps the algorithm reaches the global minimiser $w^{*}=3$ with $f(w^{*})=0$.

\newpage
\section*{Assumptions}
\begin{assumption}[Smoothness]\label{ass:smooth}
$f$ is twice continuously differentiable and its gradient is $L$‑Lipschitz:
\[
\norm{\nabla f(x)-\nabla f(y)}\le L\norm{x-y},
\qquad\forall x,y\in\R^{n}.
\]
\end{assumption}
\begin{assumption}[Bounded Hessian]\label{ass:hessian}
There exists $M>0$ such that
\[
\norm{\nabla^{2}f(x)}\le M,\qquad\forall x\in\R^{n}.
\]
\end{assumption}
\begin{assumption}[Model Accuracy]\label{ass:model}
The matrix $B_{k}$ satisfies
\[
\bigl\| (B_{k}-\nabla^{2}f(x_{k}))p\bigr\|\le \kappa\|p\|^{2},
\qquad\forall p\in\R^{n},
\]
for some constant $\kappa\ge0$ (e.g. $B_{k}= \nabla^{2}f(x_{k})$ gives $\kappa=0$).
\end{assumption}
\begin{assumption}[Algorithmic Parameters]\label{ass:params}
\[
0<\eta_{1}<\eta_{2}<1,\qquad 
\gamma_{\text{inc}}>1,\qquad 
0<\gamma_{\text{dec}}<1,\qquad 
0<\Delta_{\min}\le\Delta_{0}\le\Delta_{\max}<\infty .
\]
\end{assumption}

\section*{Main Convergence Theorem}
\begin{theorem}[Global first‑order convergence and complexity]\label{thm:main}
Assume \ref{ass:smooth}–\ref{ass:params} hold and that the Cauchy step
\eqref{1} is used as approximate solution of the trust‑region subproblem.
Let $\{x_{k}\}$ be the sequence generated by TR‑AP and let $\varepsilon>0$ be a prescribed tolerance.
Then
\begin{enumerate}[label=(\roman*)]
\item\label{it:stationary} Every limit point $\bar x$ of $\{x_{k}\}$ satisfies $\nabla f(\bar x)=0$.
\item\label{it:complexity} The algorithm produces an iterate $x_{k}$ with
\[
\norm{\nabla f(x_{k})}\le\varepsilon
\]
in at most
\[
\mathcal{O}\!\left(\frac{L\,\Delta_{\max}^{2}}{\varepsilon^{2}}\right)
\]
successful iterations. The total number of iterations (including rejected steps) is bounded by a constant multiple of the same order.
\end{enumerate}
\end{theorem}

\section*{Theoretical Properties}
\begin{itemize}
\item \textbf{Stability.} The radius update rule \eqref{3} guarantees that $\Delta_{k}$ never falls below $\Delta_{\min}>0$ and never grows beyond $\Delta_{\max}$, preventing numerical overflow or stagnation.
\item \textbf{Hyper‑parameter sensitivity.} Larger $\gamma_{\text{inc}}$ accelerates expansion when the model is accurate, but may cause occasional rejections; a smaller $\gamma_{\text{dec}}$ forces rapid contraction after a failure, improving robustness.
\item \textbf{Comparison to baseline methods.}
    \begin{itemize}
    \item Compared with plain gradient descent, TR‑AP adapts the step length via $\Delta_{k}$, yielding a guaranteed reduction even when the gradient is large.
    \item Compared with exact trust‑region solvers, the Cauchy step is cheap (\(\mathcal{O}(n)\)) while preserving the $O(\varepsilon^{-2})$ complexity.
    \end{itemize}
\end{itemize}

\section*{Discussion}
The theorem mirrors classical results for exact trust‑region methods (e.g. \cite{NocedalWright}), showing that an inexpensive approximate subproblem solution is sufficient for the same worst‑case complexity. The constants in the bound depend explicitly on the Lipschitz constant $L$ and the maximal radius $\Delta_{\max}$, which are intrinsic to the problem and the algorithmic design.

\newpage
\section*{Preliminaries (Definitions and Lemmas)}
\begin{definition}[Predicted and actual reduction]
For a given step $p$ define
\[
\text{pred}(p):=m_{k}(0)-m_{k}(p),\qquad 
\text{ared}(p):=f(x_{k})-f(x_{k}+p).
\]
\end{definition}

\begin{lemma}[Cauchy‑point reduction]\label{lem:cauchy}
Let $p_{k}$ be the Cauchy step \eqref{1}. Under Assumptions \ref{ass:smooth}–\ref{ass:hessian},
\[
\text{pred}(p_{k})\;\ge\;
\frac12\;\min\!\Bigl\{\Delta_{k},\frac{\norm{g_{k}}}{M}\Bigr\}\,\norm{g_{k}} .
\tag{4}
\]
\end{lemma}
\begin{proof}
Because $m_{k}$ is a quadratic,
\[
\text{pred}(p)= -g_{k}^{\top}p-\tfrac12 p^{\top}B_{k}p .
\]
Insert $p_{k}=-(\Delta_{k}/\norm{g_{k}})g_{k}$:
\[
\text{pred}(p_{k})=
\Delta_{k}\norm{g_{k}}-\tfrac12\Delta_{k}^{2}
\frac{g_{k}^{\top}B_{k}g_{k}}{\norm{g_{k}}^{2}} .
\]
Using $\norm{B_{k}}\le M$ (from Assumption \ref{ass:hessian}) gives
\[
\frac{g_{k}^{\top}B_{k}g_{k}}{\norm{g_{k}}^{2}}\le M .
\]
Hence
\[
\text{pred}(p_{k})\ge \Delta_{k}\norm{g_{k}}-\tfrac12 M\Delta_{k}^{2}
= \tfrac12\bigl(2\Delta_{k}\norm{g_{k}}-M\Delta_{k}^{2}\bigr).
\]
If $\Delta_{k}\le \norm{g_{k}}/M$ the term in parentheses is $\ge \Delta_{k}\norm{g_{k}}$, otherwise it is $\ge \frac{\norm{g_{k}}^{2}}{M}$. Both cases are captured by the minimum in \eqref{4}. ∎
\end{proof}

\begin{lemma}[Model–function discrepancy]\label{lem:discrep}
Under Assumption \ref{ass:smooth},
\[
\bigl|\,\text{ared}(p_{k})-\text{pred}(p_{k})\,\bigr|
\;\le\;\tfrac12 L \norm{p_{k}}^{2}.
\tag{5}
\]
\end{lemma}
\begin{proof}
By the mean‑value theorem and $L$‑Lipschitzness,
\[
f(x_{k}+p)-f(x_{k})-g_{k}^{\top}p
\;=\;\int_{0}^{1}\!\bigl[\nabla f(x_{k}+tp)-\nabla f(x_{k})\bigr]^{\!\top}p\,dt
\;\le\;\int_{0}^{1} Ltp^{\top}p\,dt
=\tfrac12 L\norm{p}^{2}.
\]
Since $m_{k}(p)-m_{k}(0)=g_{k}^{\top}p+\tfrac12 p^{\top}B_{k}p$, the $B_{k}$ term cancels in the difference, yielding \eqref{5}. ∎
\end{proof}

\section*{Main Proof (Step‑by‑Step Derivation)}
We prove Theorem \ref{thm:main}.  
All non‑trivial steps are labelled \([{\rm Step}\ N]\).

\subsection*{Step 1 – Lower bound on predicted reduction}
From Lemma~\ref{lem:cauchy} and the definition of $p_{k}$,
\[
\text{pred}_{k}\;\ge\;
\frac12\,\min\!\Bigl\{\Delta_{k},\frac{\norm{g_{k}}}{M}\Bigr\}\,\norm{g_{k}} .
\tag{6}
\]

\subsection*{Step 2 – Relating actual to predicted reduction}
Using Lemma~\ref{lem:discrep} and the definition of $\rho_{k}$,
\[
\rho_{k}
=1-\frac{\bigl|\text{ared}_{k}-\text{pred}_{k}\bigr|}{\text{pred}_{k}}
\ge 1-\frac{\tfrac12 L\norm{p_{k}}^{2}}{\text{pred}_{k}} .
\tag{7}
\]

\subsection*{Step 3 – When the step is accepted}
Assume $\rho_{k}\ge\eta_{1}$.  
From \eqref{7} we obtain a sufficient condition for acceptance:
\[
\frac{\tfrac12 L\norm{p_{k}}^{2}}{\text{pred}_{k}}\le 1-\eta_{1}.
\tag{8}
\]
Insert the bound \eqref{6} for $\text{pred}_{k}$ and $\norm{p_{k}}=\Delta_{k}$ (Cauchy step):
\[
\frac{\tfrac12 L\Delta_{k}^{2}}
{\tfrac12\min\!\{\Delta_{k},\norm{g_{k}}/M\}\,\norm{g_{k}}}
\;=\;
\frac{L\Delta_{k}}
{\min\!\{\Delta_{k},\norm{g_{k}}/M\}\,\norm{g_{k}}}
\;\le\;1-\eta_{1}.
\tag{9}
\]
Two cases arise.

\paragraph{Case A: $\Delta_{k}\le \norm{g_{k}}/M$}.  
Then $\min\{\cdot\}=\Delta_{k}$ and \eqref{9} reduces to
\[
\frac{L\Delta_{k}}{\Delta_{k}\norm{g_{k}}}= \frac{L}{\norm{g_{k}}}\le 1-\eta_{1}
\;\Longrightarrow\;
\norm{g_{k}}\ge \frac{L}{1-\eta_{1}} .
\tag{10}
\]

\paragraph{Case B: $\Delta_{k}> \norm{g_{k}}/M$}.  
Now $\min\{\cdot\}= \norm{g_{k}}/M$ and \eqref{9} becomes
\[
\frac{L\Delta_{k}}{(\norm{g_{k}}/M)\norm{g_{k}}}
= \frac{LM\Delta_{k}}{\norm{g_{k}}^{2}}
\le 1-\eta_{1}
\;\Longrightarrow\;
\norm{g_{k}}^{2}\ge \frac{LM\Delta_{k}}{1-\eta_{1}} .
\tag{11}
\]

\subsection*{Step 4 – Guaranteed decrease when $\norm{g_{k}}>\varepsilon$}
Let $\varepsilon>0$ be arbitrary. Define
\[
c_{1}:=\frac{1-\eta_{1}}{2M},\qquad
c_{2}:=\frac{1-\eta_{1}}{2L}.
\]
If $\norm{g_{k}}\ge\varepsilon$ and $\Delta_{k}\le\Delta_{\max}$, then at least one of the following holds:
\begin{enumerate}[label=(\alph*)]
\item $\Delta_{k}\le \norm{g_{k}}/M$.
\item $\Delta_{k}> \norm{g_{k}}/M$.
\end{enumerate}
Using \eqref{6} and the acceptance condition $\rho_{k}\ge\eta_{1}$ we obtain, in both cases,
\[
\text{ared}_{k}\;\ge\;\eta_{1}\,\text{pred}_{k}
\;\ge\;\eta_{1}\,\frac12\min\!\Bigl\{\Delta_{k},\frac{\norm{g_{k}}}{M}\Bigr\}\,\norm{g_{k}} .
\tag{12}
\]
Since $\min\{\Delta_{k},\norm{g_{k}}/M\}\ge
\min\bigl\{\Delta_{\min},\varepsilon/M\bigr\}$ for any successful iteration,
\[
\text{ared}_{k}\;\ge\;c\,\varepsilon^{2},
\qquad 
c:=\eta_{1}\,\min\!\Bigl\{\frac{\Delta_{\min}}{2\varepsilon},\frac{1}{2M}\Bigr\}>0 .
\tag{13}
\]

\subsection*{Step 5 – Bounding the number of successful iterations}
Let $S$ be the set of successful iterations (those with $\rho_{k}\ge\eta_{1}$). Summing \eqref{13} over $k\in S$ yields
\[
\sum_{k\in S}\text{ared}_{k}\;\ge\;c\,|S|\,\varepsilon^{2}.
\tag{14}
\]
Because $f$ is bounded below (e.g. $f_{\inf}=\inf_{x}f(x)$), the total actual reduction cannot exceed $f(x_{0})-f_{\inf}$. Hence
\[
c\,|S|\,\varepsilon^{2}\;\le\;f(x_{0})-f_{\inf}.
\tag{15}
\]
Thus
\[
|S|\;\le\;\frac{f(x_{0})-f_{\inf}}{c\,\varepsilon^{2}}
=\mathcal{O}\!\bigl(\varepsilon^{-2}\bigr).
\tag{16}
\]

\subsection*{Step 6 – Bounding rejected iterations}
A rejected step contracts the radius by the factor $\gamma_{\text{dec}}<1$. Between any two successful steps the radius can be reduced at most
\[
\frac{\log(\Delta_{\max}/\Delta_{\min})}{\log(1/\gamma_{\text{dec}})}
\]
times. Hence the total number of iterations is at most a constant multiple of $|S|$, preserving the $\mathcal{O}(\varepsilon^{-2})$ bound.

\subsection*{Step 7 – Stationarity of limit points}
Assume, for contradiction, that a subsequence $\{x_{k_{j}}\}$ converges to $\bar x$ with $\norm{\nabla f(\bar x)}>\delta>0$. Then for $j$ large enough $\norm{g_{k_{j}}}\ge\delta/2$, and by \eqref{13} each successful step yields a reduction of at least $c\delta^{2}/4$, contradicting the finiteness of $f(x_{0})-f_{\inf}$. Therefore any accumulation point must satisfy $\nabla f(\bar x)=0$.

All steps together establish statements \eqref{it:stationary} and \eqref{it:complexity} of Theorem \ref{thm:main}. ∎

\section*{Rate Derivation (With Constants)}
From \eqref{15} we can write an explicit iteration‑complexity bound:
\[
|S|\;\le\;
\frac{2\bigl(f(x_{0})-f_{\inf}\bigr)M}{\eta_{1}\,\Delta_{\min}\,\varepsilon^{2}}
\quad\text{if }\Delta_{\min}\le\frac{\varepsilon}{M},
\]
and
\[
|S|\;\le\;
\frac{2\bigl(f(x_{0})-f_{\inf}\bigr)L}{\eta_{1}\,\varepsilon^{2}}
\quad\text{otherwise}.
\]
Thus the worst‑case constant is
\[
C:=\frac{2\max\{M\Delta_{\min}^{-1},\,L\}}{\eta_{1}}\,\bigl(f(x_{0})-f_{\inf}\bigr),
\]
giving the compact bound $|S|\le C\,\varepsilon^{-2}$.

\section*{Special Cases}
\begin{enumerate}
\item \textbf{Exact Hessian ($B_{k}= \nabla^{2}f(x_{k})$).} Then $\kappa=0$ in Assumption \ref{ass:model}, and Lemma~\ref{lem:cauchy} holds with $M$ replaced by the true spectral norm of $\nabla^{2}f(x_{k})$, often yielding a larger reduction constant.
\item \textbf{Quadratic objective.} For $f(x)=\tfrac12 x^{\top}Qx+b^{\top}x+c$ the model is exact; $\rho_{k}=1$ for every step, the radius never contracts, and the algorithm terminates in at most $\lceil \norm{x_{0}-x^{*}}/\Delta_{0}\rceil$ iterations.
\item \textbf{Large initial radius.} If $\Delta_{0}\ge \norm{g_{0}}/M$, the Cauchy step is the full Newton step projected onto the trust region, and the algorithm behaves like a damped Newton method.
\end{enumerate}

\begin{thebibliography}{9}
\bibitem{NocedalWright}
J.~Nocedal and S.~J.~Wright,
\emph{Numerical Optimization},
2nd~ed., Springer, 2006.
\end{thebibliography}

\end{document}